\chapter{1. Finite propagation speed, \texorpdfstring{$L^2$}{L2} estimates, and bounded sectional curvature}
\source{cheeger-gromov-taylor:page-0005,cheeger-gromov-taylor:page-0006,cheeger-gromov-taylor:page-0007,cheeger-gromov-taylor:page-0008,cheeger-gromov-taylor:page-0009}

In this section we derive estimates on the kernel $k_f(x_1,x_2)$ for
$f(\sqrt{-\Delta})$ in an elementary fashion, under the assumption that $M$ is a
complete Riemannian manifold with bounded sectional curvature,
$\lvert K_M\rvert \leq K$.

We begin by explaining how the finite propagation speed and $L^2$ operator
norm boundedness of $\cos s\sqrt{-\Delta}$ lead to $L^2$ estimates on the
kernels of functions of $\sqrt{-\Delta}$. These estimates, which are valid for
arbitrary complete Riemannian manifolds, will also be applied in the following
two sections.

Recall that $u=\cos s\sqrt{-\Delta}\,\delta_{x_2}$ satisfies the wave equation

\begin{equation}\tag{1.1}
\left(\frac{\partial^2}{\partial s^2}-\Delta\right)u=0,
\end{equation}

with initial data

\begin{equation}\tag{1.2}
u(0,x)=\delta_{x_2}, \qquad \frac{\partial}{\partial s}u(0,x)=0.
\end{equation}

Thus by general results on hyperbolic equations (see, e.g., [32, Chapter IV]), it follows that

\begin{equation}\tag{1.3}
\supp \cos s\sqrt{-\Delta}\,\delta_{x_2}\subset \overline{B_{|s|}(x_2)},
\end{equation}

i.e., $\cos s\sqrt{-\Delta}$ has \emph{unit} propagation speed. This fact can be
brought to bear on the estimation of the kernels of other functions of
$\sqrt{-\Delta}$, by employing the formula, valid for an even function $f(\lambda)$:

\begin{equation}\tag{1.4}
f(\sqrt{-\Delta})=\frac{1}{2\pi}\int_{-\infty}^{\infty}\hat f(s)\cos s\sqrt{-\Delta}\,ds,
\end{equation}

where

\begin{equation}\tag{1.5}
\hat f(s)=\int_{-\infty}^{\infty}f(\lambda)\cos \lambda s\,d\lambda.
\end{equation}

This is just the Fourier inversion formula for even functions, granted the
spectral theorem for $\sqrt{-\Delta}$. In case $M=S^1_1$, the unit circle, (1.4)
is essentially the Poisson summation formula. This approach to the study of
$f(\sqrt{-\Delta})$ has been used by some of the authors before [32], [33], [8] in
contexts more like those considered in \S3 of this paper. The finite propagation
speed property (1.3) has entered also into other qualitative studies of
$\Delta$ on a complete Riemannian manifold. We mention particularly Chernoff's
elegant proof that $\Delta^k$ is essentially self adjoint on $C_0^\infty(M)$ for all
$k$, if $M$ is complete; see [11]. Also, Kannan [37] has obtained estimates for
the heat kernel on compact manifolds.

\begin{definition}[Finite-propagation class]
\label{def:finite-propagation-class}
\dcref{Definition 1.1}
\group{section-1}
\source{cheeger-gromov-taylor:page-0006}
We say
\begin{equation}\tag{1.6}
f \in \Tscr_1^k,
\end{equation}
provided that $f$ is even and, for each $I = (-\varepsilon, \varepsilon)$,
\begin{equation}\tag{1.7}
\left(\frac{d}{ds}\right)^l \hat{f}(s) \in L^1(\R \setminus I), \qquad 0 \leq l \leq k.
\end{equation}
\end{definition}

\begin{proposition}[Tail estimate for $f(\sqrt{-\Delta})$]
\label{prop:tail-estimate}
\dcref{Proposition 1.1}
\group{section-1}
\source{cheeger-gromov-taylor:page-0006}
\textit{Let $M^n$ be any complete Riemannian manifold, and let $x_2 \in M^n$.
Let $B_R(x_2)$ be the ball of radius $R$ about $x_2$. If $f \in \Tscr_1$, then}
\begin{equation}\tag{1.8}
\left\| f(\sqrt{-\Delta})u \right\|_{M \setminus B_R(x_2)}
\leq \|u\| \frac{1}{\pi} \int_{R-r}^{\infty} |\hat{f}(s)|\, ds
\end{equation}
\textit{provided}
\begin{equation}\tag{1.9}
\supp u \subset B_r(x_2), \qquad r < R.
\end{equation}
\end{proposition}

\begin{proof}
\source{cheeger-gromov-taylor:page-0006}
Use the formula, valid for even $f$,
\begin{equation}\tag{1.10}
f(\sqrt{-\Delta})u = \frac{1}{2\pi} \int_{-\infty}^{\infty} \hat{f}(s)\cos s\sqrt{-\Delta}\,u\,ds.
\end{equation}
Now (1.3) and (1.9) imply
\begin{equation}\tag{1.11}
\cos s\sqrt{-\Delta}\,u = 0 \quad \text{on } M \setminus B_{r+|s|}(x_2).
\end{equation}
Thus the left side of (1.8) is bounded by
\[
\left\| \frac{1}{\pi} \int_{R-r}^{\infty} \hat{f}(s)\cos s\sqrt{-\Delta}\,u\,ds \right\|
\leq \frac{1}{\pi} \int_{R-r}^{\infty} \|u\|\,|\hat{f}(s)|\,ds,
\]
since
\begin{equation}\tag{1.12}
\|\cos s\sqrt{-\Delta}\| \leq 1.
\end{equation}
This proves (1.8).
\end{proof}

\begin{corollary}[Derivative tail estimate]
\label{cor:derivative-tail-estimate}
\dcref{Corollary 1.2}
\group{section-1}
\source{cheeger-gromov-taylor:page-0006}
\textit{Under the hypotheses of Proposition 1.1, if $f \in \Tscr_1^{2k+2l}$, then}
\begin{equation}\tag{1.13}
\left\| \Delta^k f(\sqrt{-\Delta}) \Delta^l u \right\|_{M \setminus B_R(x_2)}
\leq \|u\| \cdot \frac{1}{\pi} \int_{R-r}^{\infty} \left|\hat{f}^{(2k+2l)}(s)\right|\, ds.
\end{equation}
\end{corollary}

\begin{proof}
\source{cheeger-gromov-taylor:page-0006}
The identity (1.4) shows that
\begin{equation}\tag{1.14}
\Delta^k f(\sqrt{-\Delta}) \Delta^l u
= \frac{1}{2\pi} \int_{-\infty}^{\infty} \hat{f}^{(2k+2l)}(s)\cos s\sqrt{-\Delta}\,u\,ds,
\end{equation}
so (1.13) follows exactly as (1.8).
\end{proof}

We now specialize to the case $\lvert K_M\rvert \le K$. Choose
\begin{equation}\tag{1.15}
 r_1\le \min\bigl(K^{-1/2},\ii(x_1)\bigr),\qquad
 r_2\le \min\bigl(K^{-1/2},\ii(x_2)\bigr).
\end{equation}
Then we claim that in a ball of radius $r_j$ about $x_j$, the constant in the
elliptic estimate for $\Delta$ can be controlled. Thus Corollary 1.2 leads to a
pointwise estimate in this case. To see this, first assume that
\begin{equation}\tag{1.16}
 r_j=1,\qquad \lvert K_M\rvert \leq 1,
\end{equation}
and that (1.15) holds. Recall that if $g(r)$ is any function of $r$ alone in a
normal ball $B_{r_0}(x_0)$, we have
\begin{equation}\tag{1.17}
 \Delta g = g'' + \frac{A'}{A}g'.
\end{equation}
Here the prime denotes differentiation with respect to $r$, and, in general,
$A(r,\theta)$ denotes the area element in polar coordinates, on the complement
of the cut locus of $x_0$. Thus $A/r^{n-1}$ is the Radon-Nikodym derivative of
the pull-back by $\exp^*_{x_0}$ of the volume element of $M$ with respect to
Lebesgue measure on $T_{x_0}M$. A standard comparison (see [7, Chapter 1])
shows that if $\lvert K_M\rvert \leq 1$, then on $B_1(x_0)$
\begin{equation}\tag{1.18}
 0<\frac{c_1(n)}{r}<\frac{A'}{A}<\frac{c_2(n)}{r}.
\end{equation}
Let $\phi$ be a smooth function supported on $[0,1/2]$ with
$\phi\,|_{[0,1/4]}\equiv 1$, $|\phi'|<5$, and set $\phi_{\eps}=\phi(r/\eps)$.
Let
\[
 P=\phi_{\eps}(r)\,\frac{r^{2-n}}{\alpha(n-1)},\qquad
\left(P-\phi_{\eps}\frac{\log r}{2\pi},\ n=2\right).
\]
Finally let $\Delta_x(P(x_0,x))=Q$ denote $\Delta_x$ applied pointwise to $P$.
Then on $B_1(x_0)$,
\begin{equation}\tag{1.19}
 |Q(r)|\le C_{\eps}r^{2-n}.
\end{equation}
If $g(x)$ is a smooth function on $B_1(x_0)$, then
\begin{equation}\tag{1.20}
 \int_{B_1(x_0)} P(x_0,x)\Delta g(x)\,dx = g(x_0)-\int_{B_1(x_0)} Q(x)g(x)\,dx.
\end{equation}
By the above mentioned comparison, the metric on $B_1(x_0)$ is bounded
uniformly above and below by a constant times the Euclidean metric in normal
coordinates. Fix $\eps$ sufficiently small so that the $N^2$ compositions below
are defined and apply the standard iteration argument (see e.g. [15]) to (1.20).
Thus if we write (1.20) as
\begin{equation}\tag{1.21}
 P\Delta = I - Q,
\end{equation}
and set
\begin{equation}\tag{1.22}
 \Pscr = (I + \Qscr + \cdots + \Qscr^{N-1})P,\qquad
 N=\left[\frac{n}{4}\right]+1,
\end{equation}
\[
\Qscr = Q^N,
\]
then it follows easily that
\begin{equation}\tag{1.23}
 \Pscr^N(\Delta^N g) + \Pscr^{N-1}\Qscr(\Delta^{N-1}g) + \cdots + \Pscr\Qscr(\Delta g) + \Qscr g = g.
\end{equation}
The kernels on the left hand are continuous and their sup norms are bounded,
since the metric is boundedly related to the Euclidean metric in normal
coordinates, and (1.19) holds. Thus we get
\begin{equation}\tag{1.24}
|g(x_0)| \leq c(n)\bigl(\|g\|_{B_{i(x_0)}} + \|\Delta g\|_{B_{i(x_0)}} + \cdots + \|\Delta^N g\|_{B_{i(x_0)}}\bigr),
\qquad
N=\left[\frac{n}{4}\right]+1.
\end{equation}

We can now obtain the corresponding estimate in the general case by
multiplying the metric by a suitable constant. This gives

\begin{proposition}[Local elliptic estimate]
\label{prop:local-elliptic-estimate}
\dcref{Proposition 1.3}
\group{section-1}
\source{cheeger-gromov-taylor:page-0008}
\textit{Let $\lvert K_M\rvert \leq K$.}

\begin{equation}\tag{1.25}
r \leq \min\bigl(|K|^{-1/2}, \ii(x_0)\bigr),\qquad
N=\left[\frac{n}{4}\right]+1,
\end{equation}

\textit{Then}

\begin{equation}\tag{1.26}
|g(x_0)| \leq c(n)\bigl(r^{-n/2}\|g\|_{B_r(x_0)} + \cdots + r^{2N-n/2}\|\Delta^N g\|_{B_r(x_0)}\bigr).
\end{equation}
\end{proposition}

In Theorem 4.2 we will give a lower estimate for $\ii(x)$ which will be
incorporated into Theorem 1.4 below. First we introduce some notation.
Let $\Vr(p)$ denote the volume of $B_r(p)$ in $M^n$, and $\VHr$ the volume of
$B_r(q)\subset M_H^n$, where $M_H^n$ is the simply connected $n$-dimensional
space of curvature $H$. Thus

\begin{equation}\tag{1.27}
\VHr = \alpha(n-1)\int_0^r \QH(s)\,ds,
\end{equation}

where

\begin{equation}\tag{1.28}
\QH(s)=
\begin{cases}
\left(\dfrac{\sin \sqrt{H}\,s}{\sqrt{H}}\right)^{n-1}, & H>0,\\[1.2ex]
s^{n-1}, & H=0,\\[1.2ex]
\left(\dfrac{\sinh \sqrt{|H|}\,s}{\sqrt{|H|}}\right)^{n-1}, & H<0,
\end{cases}
\end{equation}

and $\alpha(n-1)=V(S^{n-1}_1)$.

\begin{theorem}[Kernel estimates for bounded sectional curvature]
\label{thm:kernel-estimates-bounded-curvature}
\dcref{Theorem 1.4}
\group{section-1}
\source{cheeger-gromov-taylor:page-0009}
\textit{Let $M^n$ be complete and $H \leq \KM \leq K$. Fix $p$, $x_1$, $x_2 \in M^n$,
$r_0 > 0$. Set $\rho_j = \rho(p,x_j)$, $j=1,2$, and $d = \rho(x_1,x_2)$.
Let $4N > n$ and assume $\hat f^{(k)} \in L^1(\R)$, $0 \le k \le 4N$.
Fix $r$, $r_0$, $s$ with $r_0 + 2s \leq \pi/\sqrt{K}$, $r_0 \leq \pi/4\sqrt{K}$.}

\textit{(i) Then the kernel $k_f$ of $f(\sqrt{-\Delta})$ is continuous on
$\{(x_1,x_2) \in M \times M:\ x_1 \ne x_2\}$ and satisfies}

\begin{equation}\tag{1.29}
|k_f(x_1,x_2)| \leq \frac{1}{\pi} c^2(n) \sum_{0 \le i,j \le N}
\bar r_1^{2i-n/2}\bar r_2^{2j-n/2} \int_0^\infty \bigl|\hat f^{(2i+2j)}(s)\bigr|\,ds,
\end{equation}

\textit{where $c(n)$ is the constant in (1.26) and}

\begin{equation}\tag{1.30}
\bar r_j = \min\!\left(|H|^{-1/2}, \frac{r_0}{2}\,\frac{1}{1 + \left(V^{H}_{r_0+s}/V_r(p)\right)\left(V^{H}_{\rho_j+r}/V^{H}_s\right)}\right).
\end{equation}

\textit{(ii) If only $f \in \mathcal{F}^{4N}_1$, then for $r_j < \bar r_j$, $j=1,2$,
and $r_1 + r_2 < d$,}

\begin{equation}\tag{1.31}
|k_f(x_1,x_2)| \leq \frac{2}{\pi} c^2(n) \sum_{0 \le i,j \le N}
r_1^{2i-n/2}r_2^{2j-n/2} \int_{d-r_1-r_2}^\infty \bigl|\hat f^{(2i+2j)}(s)\bigr|\,ds.
\end{equation}
\end{theorem}

\begin{proof}
\source{cheeger-gromov-taylor:page-0009}
\uses{prop:local-elliptic-estimate,cor:derivative-tail-estimate}
\textit{Proof of Theorem 1.4.}\ \ Since the arguments for (1.29), (1.31) are
essentially the same, we consider (1.31). Let $u \in L^2$ be supported on
$B_{r_2}(x_2)$. We apply (1.26) of Proposition 1.3 with
$g = \Delta_2^j k_f(x_1,x_2)(u(x_2))$, $x_0 = x_1$; by (1.30) and the estimate
for $i(x_1)$ of Theorem 4.7, the hypothesis (1.25) is satisfied. If we combine
the resulting estimate with (1.13) of Corollary 1.2, we get

\begin{equation}\tag{1.32}
\bigl|\Delta_2^j k_f(x_1,x_2)(u)\bigr| \leq \frac{c(n)}{\pi}\sum_{i=0}^N r_1^{2i-n/2}
\bigl\|\Delta_1^i \Delta_2^j k_f(u)\bigr\|_{B_{r_1}(x_1)}
\end{equation}
\[
\leq c(n)\|u\|_{B_{r_2}(x_2)} \sum_{i=0}^N r_1^{2i-n/2}
\int_{d-r_1-r_2}^\infty \bigl|\hat f^{(2i+2j)}(s)\bigr|\,ds.
\]

Since $u$ is arbitrary, for all $x_1$ we have

\begin{equation}\tag{1.33}
\bigl\|\Delta_2^j k_f(x_1,x_2)\bigr\|_{B_{r_2}(x_2)}
\leq c(n)\sum_{i=0}^N r_1^{2i-n/2} \int_{d-r_1-r_2}^\infty \bigl|\hat f^{(2i+2j)}(s)\bigr|\,ds.
\end{equation}

If we now apply (1.26) on $B_{r_2}(x_2)$, (1.31) follows.
\end{proof}

\begin{remark}
\label{rem:sectional-curvature-kernel-estimate}
\dcref{Remark after Theorem 1.4}
\group{section-1}
\source{cheeger-gromov-taylor:page-0009}
The estimate for $i(x_j)$ incorporated in Theorem 1.4 can be sharpened
somewhat further, if $\rho_j$ is sufficiently large; see Theorem 4.7(ii). Also,
one can easily choose the constants $r$, $r_0$, $s$, so that the estimate is
optimal. In particular cases in which $M^n$ is compact, one can make use of
known estimates on the injectivity radius; see e.g. [7].
\end{remark}

We point out that Theorem 1.4 can be extended to the case $\partial M^n \neq \emptyset$, by
use of a suitable generalization of Proposition 4.2.
